\begin{center}
    \large \textbf{Résumé : Masked Autoencoders Are Scalable Vision Learners} \\
    \vspace{0.5cm}
    Cédric GAUTHERET, Romain AMIGON et Noé BARRANGER
\end{center}

Les auteurs dans \footcite{mae22} s'attaquent à la problématique du pré-entraînement auto-supervisé en vision par ordinateur, motivés par l'écart de performance avec le NLP, où des méthodes comme BERT ont démontré la puissance du masked autoencoding. Pour ce faire, ils proposent une approche nommée Masked Autoencoder (MAE), reposant sur deux choix de conception centraux. Premièrement, une architecture encodeur-décodeur asymétrique : l'encodeur, un Vision Transformer (ViT), traite uniquement les patches visibles de l'image (sans tokens de masquage), tandis qu'un décodeur léger reconstruit l'image complète à partir des représentations latentes et des tokens de masquage réintroduits à ce stade. Deuxièmement, un taux de masquage élevé (75\%), bien supérieur aux 15\% utilisés par BERT, justifié par la redondance spatiale des images : un masquage faible permettrait une reconstruction par simple interpolation des patches voisins, sans nécessiter de compréhension globale. La cible de reconstruction est constituée des valeurs de pixels (normalisées par patch) des zones masquées, avec une perte MSE calculée uniquement sur ces zones.

Cette conception asymétrique réduit considérablement le coût de calcul puisque l'encodeur ne traite que 25\% des patches, permettant un gain de vitesse d'entraînement de 3$\times$ et plus, ainsi qu'une réduction de la consommation mémoire facilitant le passage à l'échelle vers des modèles de grande capacité. Les auteurs valident leur approche sur ImageNet-1K en pré-entraînement, puis évaluent par fine-tuning et linear probing. Avec un ViT-Huge, MAE atteint 87.8\% de précision top-1 sur ImageNet-1K (résolution 448), dépassant l'état de l'art parmi les méthodes utilisant uniquement IN1K. Le transfer learning est également démontré sur la détection et la segmentation d'objets sur COCO (53.3 AP\textsubscript{box} avec ViT-L, contre 49.3 pour le pré-entraînement supervisé) ainsi que sur la segmentation sémantique sur ADE20K (53.6 mIoU, contre 49.9 en supervisé).

Les ablations menées par les auteurs apportent plusieurs résultats notables : un encodeur sans tokens de masquage améliore la précision de 14\% en linear probing tout en accélérant l'entraînement, la profondeur du décodeur influence significativement la qualité des représentations latentes (un décodeur suffisamment profond améliore le linear probing de 8\%), et le pré-entraînement bénéficie de schedules longs (1600 epochs sans saturation observée). Les auteurs comparent également différentes stratégies de masquage (aléatoire, par blocs, par grille) ainsi que différentes cibles de reconstruction (pixels, pixels normalisés, tokens dVAE), montrant que les choix simples (masquage aléatoire et pixels normalisés) sont les plus efficaces. MAE se distingue ainsi par sa simplicité conceptuelle, sa scalabilité et son efficacité computationnelle, bien qu'il produise des représentations moins linéairement séparables que les méthodes contrastives, nécessitant un fine-tuning pour être pleinement exploité.